% New J/M extension; prerequisite definitions are in es_niemeier_cubic_arithmetic.tex.
\section{The cubic on the common lattice and the fourth Niemeier extension}
\label{jmv:sec:JM-extension}

Retain the original coordinates, Cayley multiplication, $F$, the generators $g_i$
and their heights $H_i$ from Section~\ref{nca:sec:cubic-arithmetic}, and the
exact common lattice and fourth extension of Section~\ref{esrf:sec:fixed-neighbor}.
The symbol $u$ below is the unchanged vector $u_*$ in (RF32).
\subsection{Integral generators of the common lattice and the fixed extension}

Define the unchanged residue action by
\[
 (rX)_{i,j}=x_{i,j-2\pmod6},\qquad (rX)_{D,k}=d_k.
\]
The two primary parts of the original glue have the exact descriptions
\begin{align*}
 H_2={}&\{(\varepsilon_1+\varepsilon_2+\varepsilon_3,
 \varepsilon_1+\varepsilon_2,\varepsilon_1+\varepsilon_3,
 \varepsilon_1;\varepsilon_2,\varepsilon_3):\varepsilon_i\in\mathbb F_2\},\\
 H_3={}&\{(t+s,t+2s,t,s;00):t,s\in\mathbb F_3\}.
\end{align*}
In the original \((\mathbb Z/6)^4\) coordinates this means
\(k=3w+2z\), with \(w\) the four binary entries of \(H_2\)
and \(z\) the four ternary entries of \(H_3\). The binary
\(D_4\) marking is retained. Chinese remainders recover \(w,z\)
uniquely, and the displayed parameterizations are injective.
Consequently
\[
 \beta:N\longrightarrow\mathbb F_3,\qquad \beta(X)=z_4=s
\]
is a well-defined additive homomorphism. The actual lifts
\eqref{nca:eq:N-generators} satisfy
\[
 \beta(g_i)=0\quad(1\leq i\leq28),\qquad \beta(g_{29})=1.
\]
Indeed roots have zero glue, the first three lifts are two-primary,
and the last two have \((t,s)=(1,0),(0,1)\), respectively.

Set, with no change of representatives,
\begin{align}
 I&=\{X\in N:(\rho,X)\equiv0\pmod6,\ \beta(X)=0\},
       &J&=I+\mathbb Z(3v),\label{jmv:eq:I-J}\\
 \omega_t&=(\underbrace{1,\ldots,1}_{t},0,\ldots,0)
                    -\frac t6(1,1,1,1,1,1),
       &u&=(\omega_2,\omega_4,0,\omega_2;0),\notag\\
 p&=u-r(u),&q&=v+r(u)-r(v),\notag\\
 M&=J+\mathbb Zp+\mathbb Zq.&&\label{jmv:eq:M}
\end{align}
Thus \(J\) and \(M\) are precisely the common-lattice and
residue-fixed extension presentations being evaluated here; the
following value-group proof needs no root count, lattice determinant,
or classification label for \(M\). For explicitness, the two extra
vectors in the original five blocks are
\begin{align*}
 p={}&((1,1,-1,-1,0,0),(1,1,0,0,-1,-1),0,
                         (1,1,-1,-1,0,0);0),\\
 q={}&\tfrac13((-2,4,4,-2,-2,-2),0,(2,2,-1,-1,-1,-1),
                         (1,1,1,1,-2,-2);0).
\end{align*}
The zero blocks have their original sizes six and four. These
coordinates follow by substituting \(v=\rho/6-\alpha\) into
\eqref{jmv:eq:M}; in particular the shifted \(\alpha\) term in
\(r(v)\) is retained. Also \(u\in N\),
\(\beta(u)=1\), \((\rho,u)=12\), and \((u,u)=4\):
its class is \((2,4,0,2;00)\),
\((\rho_A,\omega_t)=t(6-t)/2\), and
\((\omega_t,\omega_t)=t(6-t)/6\).

\begin{theorem}[Integral generating lists with both kernel conditions]
\label{jmv:thm:IJM-generators}
Define the following 32 ordered vectors:
\begin{align}
 b_1&=6\alpha,&
 b_i&=g_i-H_i\alpha\quad(2\leq i\leq28),\notag\\
 b_{29}&=3g_{29}-60\alpha,&
 b_{30}&=3v=\rho/2-3\alpha,\notag\\
 b_{31}&=p,&b_{32}&=q.\label{jmv:eq:IJM-generators}
\end{align}
Then the integral, not merely rational, equalities are
\[
 I=\sum_{i=1}^{29}\mathbb Zb_i,\qquad
 J=\sum_{i=1}^{30}\mathbb Zb_i,\qquad
 M=\sum_{i=1}^{32}\mathbb Zb_i.
\]
\end{theorem}
\begin{proof}
Every one of the first 29 vectors has \(\beta=0\) and height
divisible by six, by the exact heights and characters already
calculated. Conversely take any integral expression
\(X=\sum_{i=1}^{29}n_ig_i\) for \(X\in I\). Applying
\(\beta\) shows \(n_{29}=3a\) for an integer \(a\),
regardless of which expression was chosen. The exact equality
\[
 X=\sum_{i=2}^{28}n_i(g_i-H_i\alpha)
        +a(3g_{29}-60\alpha)
        +\left(n_1+\sum_{i=2}^{28}n_iH_i+60a\right)\alpha
\]
has last coefficient \((\rho,X)\), which is divisible by six.
It therefore expresses \(X\) as an integer combination of the
first 29 \(b_i\). The reverse containment was proved for each
generator. The last two equalities follow by adjoining exactly
the vectors in \eqref{jmv:eq:I-J} and \eqref{jmv:eq:M}.
No assertion of linear independence is used.
\end{proof}

\subsection{Complete finite certificates for the two new value groups}

\begin{theorem}[Equality of the common-lattice and fixed-extension value-generated groups]
\label{jmv:thm:JM-values}
For the retained Cayley multiplication and the original ordered
coordinates, the two additive cubic-value groups are exactly
\begin{equation}
 \mathcal V_F(J)=\mathcal V_F(M)
 =\frac1{27}\mathbb Z\ \oplus\frac{\sqrt2}{216}\mathbb Z
                         \ \oplus\frac{\sqrt3}{18}\mathbb Z.
 \label{jmv:eq:JM-values}
\end{equation}
The equality concerns generated additive groups, not equality
of individual-value sets \(F(J)\) and \(F(M)\).
\end{theorem}
\begin{proof}
Apply Theorem~\ref{nca:thm:finite-values} to the first 30 and
all 32 vectors in \eqref{jmv:eq:IJM-generators}. The resulting
list lengths are exactly
\[
 30+2\binom{30}{2}+\binom{30}{3}=4960,\qquad
 32+2\binom{32}{2}+\binom{32}{3}=5984.
\]
Write all pair and triple symbols in the next tables using these
\(b_i\), not the \(g_i\) or \(h_i\). Substitution in the full
polynomial gives the complete denominator certificate
\begin{equation}
\begin{array}{c|c|c|c|c}
 &A_i&D^+_{ij}&D^-_{ij}&D_{ijk}\\\hline
 J&(27,216,1)&(9,36,18)&(9,36,18)&(9,36,18)\\
 M&(27,216,1)&(9,36,18)&(9,36,18)&(9,36,18)
\end{array}.
\label{jmv:eq:JM-denominators}
\end{equation}
As before each triple is the least common multiple of all reduced
rational denominators in its full category, in the order
\((P_0,P_2,P_3)\); a zero has denominator one. The finite
arithmetic is completely specified by
\eqref{nca:eq:original-s}--\eqref{nca:eq:P3},
\eqref{jmv:eq:IJM-generators}, and all indices in
\eqref{nca:eq:single}--\eqref{nca:eq:Dthree}.
The accompanying standard-library checker
\texttt{check\_jm\_values.py}, using the included
\texttt{cubic\_engine.py}, executes every entry with exact
fractions and asserts this table. It reconstructs all 72 marked
glue classes, checks all 29 \(\beta(g_i)\), and checks both
kernel conditions on all the displayed \(I\) generators.
These finite computations establish the upper containment in
\eqref{jmv:eq:JM-values}, including all mixed polarization terms
involving \(p,q\), not just their individual cubic values.

Here is a short certificate for the reverse containment, using
only \(J\). Coordinates of each row are taken in the ordered
proposed value-group basis
\((1/27,\sqrt2/216,\sqrt3/18)\):
\begin{equation}
\begin{array}{c|r r r}
 F(b_2)&14&-272&0\\
 F(b_4)&5&-128&0\\
 F(b_{27})&23220&-111672&-90\\
 F(b_{28})&13780&-74672&0\\
 F(b_{30})&378&-1701&-18\\
 D^+_{2,27}&7425&-60624&-3\\
 D^+_{11,21}&36&48&-16
\end{array}.
\label{jmv:eq:JM-lower-certificate}
\end{equation}
Let \(C_1,C_2,C_3\) have, respectively, columns obtained by
transposing rows \((2,3,5)\), \((2,5,6)\), and \((1,4,7)\)
of this table. Direct three-by-three determinants are
\begin{align*}
 d_1&=-39859290,&d_2&=11531403,&d_3&=-43244032.
\end{align*}
The following explicit integer identity certifies their gcd one:
\begin{equation}
 (-6882196252)d_1+(-23788908901)d_2+(-907)d_3=1.
 \label{jmv:eq:JM-bezout}
\end{equation}
Put \((c_1,c_2,c_3)=(-6882196252,-23788908901,-907)\).
For each standard basis column \(e_t\in\mathbb Z^3\) the exact
integer reconstruction is
\begin{equation}
 e_t=\sum_{j=1}^3 C_j\bigl(c_j\operatorname{adj}(C_j)e_t\bigr).
 \label{jmv:eq:JM-axis-reconstruction}
\end{equation}
All columns of every \(C_j\) represent integer combinations of
actual cubic values on \(J\); their integer coefficient vectors
in \eqref{jmv:eq:JM-axis-reconstruction} therefore give each
of the three claimed generators as such a combination. This is
an explicit reverse-containment morphism, not an index inference.
The checker also expands every adjugate and verifies all three
axis reconstructions, saving their integer coefficients in
\texttt{EXACT\_RESULTS.json}. Thus \(\mathcal V_F(J)\) contains
the right-hand side of \eqref{jmv:eq:JM-values}. Since \(J\subset M\),
the same certificate proves its containment in \(\mathcal V_F(M)\).
The upper bounds already established finish both equalities.
\end{proof}

For additional exact checks on the unchanged representatives, the
three individual values are
\[
 F(3v)=14-\frac{63}{8}\sqrt2-\sqrt3,\qquad
 F(p)=-2,\qquad F(q)=\frac{64}{27}.
\]
They agree with the corresponding rows of the complete calculation;
they are not substituted for the mixed-term certificate.

\subsection{Exact inclusions, quotient maps, and the factor of the Albert determinant}

\begin{theorem}[The exact four-to-three chain of value groups]
\label{jmv:thm:chain}
There is a strict chain of additive subgroups of the original real
coefficient field
\[
 \mathcal V_F(N)\ \subsetneq\ \mathcal V_F(J)=\mathcal V_F(M)
                       \ \subsetneq\ \mathcal V_F(\Lambda).
\]
The inclusion matrices in the ordered bases already displayed are
\[
 \mathcal V_F(N)\longrightarrow\mathcal V_F(J)=\mathcal V_F(M):
       (a,b,c)\longmapsto(a,4b,c),
\]
\[
 \mathcal V_F(J)=\mathcal V_F(M)\longrightarrow\mathcal V_F(\Lambda):
       (a,b,c)\longmapsto(a,b,3c).
\]
Their exact quotient maps are
\begin{align*}
 \frac a{27}+\frac{b\sqrt2}{216}+\frac{c\sqrt3}{18}
       &\longmapsto b\pmod4,\\
 \frac a{27}+\frac{b\sqrt2}{216}+\frac{c\sqrt3}{54}
       &\longmapsto c\pmod3.
\end{align*}
In particular
\begin{align*}
 \mathcal V_F(J)/\mathcal V_F(N)&\cong\mathbb Z/4,&
 \mathcal V_F(M)/\mathcal V_F(N)&\cong\mathbb Z/4,\\
 \mathcal V_F(\Lambda)/\mathcal V_F(J)&\cong\mathbb Z/3,&
 \mathcal V_F(\Lambda)/\mathcal V_F(M)&\cong\mathbb Z/3,\\
 \mathcal V_F(M)/\mathcal V_F(J)&=0.&&
\end{align*}
\end{theorem}
\begin{proof}
The first basis comparison uses
\(\sqrt2/54=4(\sqrt2/216)\), with the other basis entries equal.
The second uses \(\sqrt3/18=3(\sqrt3/54)\), with its other
entries equal. Uniqueness of rational radical coefficients proves
that the displayed kernels are exactly the corresponding subgroups.
Surjectivity is witnessed by \(\sqrt2/216\) and \(\sqrt3/54\),
respectively, both actual elements of the proved additive groups.
Their orders modulo the smaller groups are exactly four and three:
the coefficient divisibility conditions are necessary and sufficient.
This also proves strictness. The last quotient is zero by the equality
of the two proved groups. In particular the composite inclusion is
\((a,b,c)\mapsto(a,4b,3c)\), recovering the full quotient
\(\mathbb Z/4\oplus\mathbb Z/3\), with explicit map
\[
 \frac a{27}+\frac{b\sqrt2}{216}+\frac{c\sqrt3}{54}
       \longmapsto (b\bmod4,c\bmod3).
\]
Its identification with \(\mathbb Z/12\) is
\((\bar b,\bar c)\mapsto9b+4c\pmod{12}\), whose inverse is
reduction modulo four and modulo three. Every assertion follows
from the proved cubic coefficient groups, not from an index of
the original lattices.
\end{proof}

The factor two in the original Albert determinant is unchanged:
\[
 \det Q_{\mathbb O}(\lambda;\mathcal LX)
       =\lambda^3-\lambda(X,X)+2F(X).
\]
Thus, if one instead asks for the additive group generated by the
actual cubic term \(2F\), it is exactly \(2\mathcal V_F(L)\).
For \(J,M\) this is
\[
 \frac2{27}\mathbb Z\ \oplus\frac{\sqrt2}{108}\mathbb Z
                           \ \oplus\frac{\sqrt3}{9}\mathbb Z.
\]
Multiplication by two is an injective additive map of the real
coefficient field and restricts to the exact isomorphism
$\mathcal V_F(L)\to\mathcal V_{2F}(L)=2\mathcal V_F(L)$,
$x\mapsto2x$, with inverse $y\mapsto y/2$.
It intertwines the same two inclusion matrices and
the same quotients; no factor has been absorbed into \(F\).

\subsection{Reproduction and scope of the exact certificate}

The accompanying evidence folder is
\begin{flushleft}\small
\texttt{supporting\_materials/workbench/evidence/}\\
\texttt{period\_lattice\_continuation\_20260906/jm\_cubic/}
\end{flushleft}
It contains both full source files \texttt{cubic\_engine.py} and
\texttt{check\_jm\_values.py}. Running the latter with Python's
standard library performs the following finite proof certificate:
all \(20399\) polarization entries for \(N,\Lambda,J,M\), all
displayed denominator and determinant assertions, the integer
Bezout and adjugate axis reconstructions, and the exact inclusion
coefficient comparisons. The output \texttt{EXACT\_RESULTS.json}
retains each selected column, determinant, Bezout coefficient,
and every expanded integer linear combination reconstructing a
value-group axis.

There is additionally an independent coefficient-identity check:
the original 24 simple roots form a real basis of the original
24-dimensional space. The direct Cayley--Dickson implementation
and the rational polynomial implementation are compared on the
24 basis vectors, their 552 ordered-by-index plus/minus pairs,
and their 2024 three-element sums, a total of 2600 exact inputs.
The finite polarization proof shows these evaluations determine
every coefficient of a homogeneous cubic in those 24 variables.
Consequently this is a full polynomial-identity certificate on
the real ambient space, not sampling of lattice points. The
field implementation retains all four rational coefficients in
\(\mathbb Q(\sqrt2,\sqrt3)\), and verifies that the \(\sqrt6\)
coefficient is identically zero. A further 65 direct checks cover
the displayed new generating lists and \(u,p,q\).

This computation proves the stated additive value groups and
their morphisms. It asserts neither equality of individual-value
sets nor that the cubic group reconstructs a lattice or classifies
all cubic-preserving lattice isometries.
